\section{Limitations and Future Work}\label{sec:limits}

\sys{} is intentionally lightweight and non-invasive, which brings concrete boundaries that
we state explicitly below.

\paragraph{Decode-centric byte model.}
\sys{} estimates each weight-streaming launch's traffic as a single read of its referenced
weight tensor (plus same-module scale tensors), ignoring activation and cache traffic.
This approximation holds for low-batch decode, where weights dominate memory movement,
and it successfully flags launches whose effective bandwidth falls below the measured
ceiling.
It does not apply to prefill or large-batch decode, where activations and compute dominate,
nor does it directly measure the write traffic of non-weight kernels---in our case study,
\sys{} localized the anomaly to GDN \texttt{out\_proj} and identified the preceding GDN
decode kernel, while the exact write volume and dirty-L2 mechanism were established by
controlled isolation experiments and kernel inspection.

\paragraph{Pointer-scanning attribution and opaque launches.}
Parameter resolution scans 8-byte words in captured kernel argument buffers for addresses
falling within the exported memory map.
Kernels that receive tensor addresses indirectly---for example, through device-side pointer
arrays or indirection tables---cannot be attributed to specific weights.
For direct runtime-API launches, \sys{} resolves parameter layouts either through the
application's loaded \texttt{cudart} or via a by-name fallback table populated from both
instantiated CUDA graph nodes and direct launches that expose a \texttt{CUfunction} handle.
However, some runtime launches arrive at the CUPTI callback with neither a function handle
nor a symbol name, and neither the cudart lookup nor \texttt{dladdr} on the host stub
resolves them, so neither path applies. We suspect a cudart other than the one \sys{} queries
(for example, a statically linked copy), but have not confirmed the cause.
In our batch-1 GB10 window (38,880 launches across 40 steps, or 972 per step), all 927
graph kernel nodes per step resolve, while 22 of the 45 direct launches per step (880 in
the window) carry no parameters (\cref{sec:coverage}).
One of those 22 is a draft-step \texttt{lm\_head} launch of the same CUTLASS kernel that
already appears in the captured graphs; its layout is known, and it is dropped only because
its name is absent at launch time.

\paragraph{Experimental scope and methodology.}
Our evaluation covers one hardware architecture (NVIDIA GB10, \texttt{sm\_121}), one model
family (Qwen3.8-27B NVFP4 with \texttt{qwen3\_5\_mtp}, $k{=}3$), and one serving engine
(vLLM 0.31.0, 0.34.0, and \texttt{main}).
Because GB10 sits at its software power cap, whole-step wall-clock drift is comparable to
the batch-1 effect size, making per-kernel and GPU-busy measurements essential for
batch-1 comparisons.
We also note two caveats in the case study itself: first, two runs fell back from CUDA
graphs to eager execution and were excluded---the second batch-1 run of the last-position-only probe and one
unmodified batch-4 baseline run. We did not determine the cause; since an unmodified run fell
back as well, we do not attribute it to the probe's corrupted state; second, our deferred-store
CUDA kernel inherits the assumption of upstream vLLM PR~\#61434 that a request with
recorded speculative inputs will replay them on a subsequent speculative step, leaving a
transition to non-speculative decode untested.
Finally, in the interest of accurate provenance, the dirty-L2 mechanism on GB10 was first
diagnosed manually through trace inspection and isolated benchmarks; \sys{} was built
afterward to systematize the workflow and automatically reproduces the finding without
manual intervention.

\paragraph{Future work.}
Two extensions would remove the main manual steps in our current workflow.
First, pairing \sys{} with per-launch read and write sets collected via dynamic binary
instrumentation such as NVBit~\cite{villa2019nvbit} would quantify dirty write volumes
directly and detect dead writes---addresses written by one launch and overwritten before any
subsequent launch reads them---without model-specific knowledge.
Second, extending the accounting model to multi-GPU serving would track collective
communication launches (such as AllReduce and AllGather) alongside weight streams and cache
traffic, enabling the same invariant-based attribution on tensor- and pipeline-parallel
deployments.
